\section{3DGS による視点生成と教師生成}
\label{sec:view_generation}

アライメントが受理されたら、
そのシーンは校正済みの画像生成器として使える。
本節では、どの視点を生成し、
同じ変換からどの教師を作るかを述べる。

\begin{figure}[t]
  \centering
  \resizebox{\linewidth}{!}{% 視点生成と教師生成。すべて本稿のために作成した TikZ 図である。
\begin{tikzpicture}[
  font=\sffamily\scriptsize,
  panel/.style={draw=black!28, rounded corners=2.5pt, fill=none,
    inner sep=3mm, line width=0.45pt},
  title/.style={font=\sffamily\bfseries\footnotesize, text=black!85},
  tinylabel/.style={font=\sffamily\fontsize{6.3}{7.2}\selectfont,
    text=black!70, align=center},
  cam/.style={regular polygon, regular polygon sides=3, minimum size=2.6mm,
    inner sep=0pt, draw=poseorange, fill=poseorange!25, line width=0.45pt},
  ptdot/.style={circle, inner sep=0pt, minimum size=1.3mm, fill=gsblue},
  arr/.style={-{Stealth[length=1.7mm,width=1.2mm]}, line width=0.65pt},
  hull/.style={draw=gsblue, dashed, line width=0.6pt},
]
  % ---------- (a) SfM 制約付き軌道候補 ----------
  \begin{scope}[local bounding box=pa]
    \coordinate (o) at (1.75,1.55);
    \draw[hull, fill=gsblue!5] (0.05,0.90) -- (1.35,0.65) -- (3.45,1.05)
      -- (3.30,2.35) -- (0.70,2.50) -- cycle;
    \foreach \a in {20,70,140,215,290} {
      \node[ptdot] at ($(o)+(\a:1.45)$) {};
    }
    \draw[courtgreen, line width=0.7pt] ($(o)+(1.05,0)$) arc (0:345:1.05);
    \draw[courtgreen, line width=0.7pt] (1.15,1.25) -- (2.35,1.25)
      (1.15,1.25) -- (1.15,1.90) (2.35,1.25) -- (2.35,1.90)
      (1.15,1.90) -- (2.35,1.90);
    \draw[courtgreen, line width=0.4pt] (1.75,1.25) -- (1.75,1.90);
    \node[cam] at ($(o)+(40:1.05)$) {};
    \node[cam] at ($(o)+(150:1.05)$) {};
    \node[cam] at ($(o)+(268:1.05)$) {};
    \draw[poseorange, line width=0.45pt]
      ($(o)+(40:1.05)$) -- (1.95,1.55);
    \draw[poseorange, line width=0.45pt, dashed]
      ($(o)+(150:1.05)$) -- (1.55,1.55);
    \node[tinylabel, anchor=north, text width=44mm] at (1.75,0.35)
      {点 = 復元カメラ、実線 = 軌道候補。\\破線は凸包の内側オフセット境界};
  \end{scope}

  % ---------- (b) 弧長一様サンプリングと分割 ----------
  \begin{scope}[shift={(4.7,0)}, local bounding box=pb]
    \draw[black!22, line width=0.4pt] (0,0) rectangle (3.3,2.2);
    \draw[poseorange, line width=0.7pt] (0.35,1.1) .. controls (1.4,2.05) and (2.1,0.25) .. (3.0,1.25);
    \foreach \x/\y in {0.35/1.10, 0.72/1.44, 1.10/1.68, 1.55/1.72, 2.00/1.55,
                        2.42/1.28, 2.75/1.24, 3.00/1.25} {
      \node[cam] at (\x,\y) {};
    }
    \node[tinylabel, anchor=north west] at (0.05,2.14)
      {3D 弧長で一様に並べる};
    \node[tinylabel, anchor=north, text width=44mm] at (1.65,-0.05)
      {近接フレームを跨がせないよう、
        trajectory group 単位で train/val/test を分ける};
    \node[tinylabel, text=gsblue, anchor=north west] at (0.08,2.12)
      {円 / 楕円 / 矩形 / 超楕円、平面または正弦高さ};
  \end{scope}

  % ---------- (c) 同一変換からの整合教師 ----------
  \begin{scope}[shift={(9.6,0)}, local bounding box=pc]
    \draw[fill=softgray!75, draw=black!25, line width=0.35pt]
      (0.35,1.00) -- (3.10,1.00) -- (2.60,1.95) -- (0.85,1.95) -- cycle;
    \draw[courtgreen, line width=0.5pt] (0.85,1.95) -- (2.60,1.95)
      (0.73,1.48) -- (2.93,1.48)
      (1.20,1.00) -- (1.07,1.95)
      (2.35,1.00) -- (2.40,1.95);
    \node[ptdot] at (1.72,1.95) {};
    \node[ptdot] at (0.97,1.95) {};
    \node[ptdot] at (2.49,1.95) {};
    \node[ptdot] at (0.80,1.48) {};
    \node[ptdot] at (2.87,1.48) {};
    \node[tinylabel, text=gsblue, anchor=south] at (1.72,2.05)
      {\KPName{} + 領域 + 二値線 + 線の意味};
    \draw[poseorange, line width=0.6pt, ->] (3.70,1.40) -- (2.20,1.40);
    \node[cam] at (3.90,1.40) {};
    \node[tinylabel, text=poseorange, anchor=south] at (3.90,1.58)
      {カメラ姿勢\\と焦点距離};
    \node[tinylabel, anchor=north, text width=50mm] at (1.60,0.92)
      {同一画像フレーム内で全教師の可視性と座標が一致する};
  \end{scope}

  \node[panel, fit=(pa)] (fa) {};
  \node[panel, fit=(pb)] (fb) {};
  \node[panel, fit=(pc)] (fc) {};
  \node[title, anchor=south] at ($(fa.north)+(0,1.3mm)$) {(a) 軌道候補};
  \node[title, anchor=south] at ($(fb.north)+(0,1.3mm)$) {(b) 弧長一様サンプルと分割};
  \node[title, anchor=south] at ($(fc.north)+(0,1.3mm)$) {(c) 整合教師の生成};

  \draw[arr] (fa.east) -- (fb.west);
  \draw[arr] (fb.east) -- (fc.west);
\end{tikzpicture}
}
  \caption{\textbf{視点生成と教師生成。}
  軌道候補は復元カメラの水平凸包を内側へオフセットした領域へ解析的に制限する。
  経路は三次元弧長で一様にサンプルし、
  軌道グループ単位で train/validation/test を分ける。
  各レンダ結果には、同じ変換から
  キーポイント、領域、二値線、線の意味分類、カメラ姿勢が同時に付く。
  }
  \label{fig:view_generation}
\end{figure}

\subsection{SfM 制約付き軌道}
\label{sec:trajectories}

3DGS は観測されたカメラの近傍では忠実だが、
撮影範囲の外側では品質を保証しない。
したがって視点を広げる操作は、
再構成が支えられる範囲へ明示的に制限する必要がある。

本手法は会場中心と各コート中心に対し、
基準半径を復元カメラの水平オフセットの 90 パーセンタイルから取る。
これは提案尺度を会場ごとに適応させる。
型付き直積が形状（円・楕円・矩形・超楕円）、
半径スケール、軸比、面内方位、基準高さ、鉛直変動を列挙する。
局所経路は
\begin{align}
  \bm{\gamma}(\theta)_{xy}
    &=\mathbf{R}(\phi)
      \begin{bmatrix}a\cos\theta\\b\sin\theta\end{bmatrix},
      &a&=r_{90}\alpha, &b&=a\beta, \\
  \gamma_z(\theta)
    &=h_0+A\sin(n\theta+\varphi)
  \label{eq:trajectory}
\end{align}
の形を取る。\(\beta=1\) が円、\(A=0\) が平面経路である。
すべての高さは正でなければならない。

\paragraph{凸包による解析的制限}
候補の採否はサンプル点ではなく曲線全体で判定する。
各中心について復元カメラの水平座標の凸包を取り、
その頂点平均を基準に \(1+\epsilon\) 倍へ広げ、
さらに余白 \(0.5\) m だけ内側へオフセットする。
単位尺度の形状写像 \(\mathbf{A}\) に対し、
半平面法線 \(\mathbf{n}\) 方向の厳密な支持関数を
\(\mathbf{A}^{\top}\mathbf{n}\) で評価できるため、
「サンプルした頂点が入っている」ではなく
\emph{軌道全体が入っている}ことを線形制約として保証できる。
中心が有効領域の外ならエラーにする。
この制約はカメラ位置のヒューリスティックであり、
画質の保証ではない。この限界は\cref{sec:discussion} で述べる。

\paragraph{弧長一様サンプリング}
楕円上や鉛直変動のある経路上では、
\(\theta\) の等間隔は距離の等間隔ではない。
パラメータを 4096 点でサンプルし、
三次元の累積弧長関数を作り、
弧長で一様に補間する。
閉ループの隣接間隔が設定上限（既定 1.05 m）以下になるまで点数を増やす。
候補順序と選択は固定 seed と明示的な安定フィールド順序を使う。

\paragraph{軌道グループ単位の分割}
経路全体は分割の前にサンプルし尽くす。
同じ trajectory group に属するフレームは
必ず同じ train / validation / test のいずれかに入る。
これにより、隣接するほぼ同一の視点が分割を跨いで漏洩することを防ぐ。
分割比率は 8:1:1 である。

ただし本研究は、この分割が
\emph{未知会場への汎化}を保証すると主張しない。
軌道グループ分離が支えるのは
「同じ会場の未学習視点」の一部であり、
未知会場はシーン集合そのものを分けたときに初めて評価できる。

\begin{table}[t]
  \centering
  \caption{\textbf{軌道とサンプリングの方針。}
  数値は設定上の下限または予算であり、
  実際に採用された視点分布の証明ではない。
  採用後の分布は manifest から集計して報告する。
  }
  \label{tab:data_policy}
  \small
  \begin{tabularx}{\linewidth}{@{}p{0.26\linewidth}X@{}}
    \toprule
    次元 & 設定値 \\
    \midrule
    平面経路 & 円・楕円・矩形・超楕円。軸比 \(1.0,0.8,0.65\)。
      方位 \(0^\circ,45^\circ,90^\circ\) \\
    半径と高さ & 復元オフセットの \(0.65\)--\(0.95\) 倍。
      基準高さ \(1.5,2.25,3.0\) m \\
    鉛直経路 & 平面または正弦。振幅 \(0.0,0.25,0.5\) m \\
    視野 & 注視点高さ \(0\) または \(1.5\) m。
      水平画角 \(45^\circ\)--\(110^\circ\)（会場別の設定） \\
    サンプリング & seed 695。三次元弧長一様。隣接間隔の上限 1.05 m \\
    被覆 & 提案予算 4800。軌道グループ下限 24、採用フレーム下限 2000 \\
    分割と描画 & 軌道グループ単位で 8:1:1。8 shard \\
    SfM 制約 & 凸包を 0.5 m 内側へオフセット。空間被覆セル 1.0 m \\
    \bottomrule
  \end{tabularx}
\end{table}


\subsection{レンダリングと教師}
\label{sec:labels}

メートル系カメラは正規化シーン単位へ変換され、
公開されたレンダラへ渡される。
レンダラは RGB、深度、alpha、画素支持の証拠を返す。
本研究のパイプラインは、
別途描画したコート画像を貼り付けることをしない。
したがって背景、面の外観、遮蔽の構造は、
元動画から再構成されたままである。

正準コートの物理点 \(\mathbf{X}_k\) について、
\begin{equation}
  \mathbf{X}^{\Cam}_k =
  \T{\Cam}{\Scene_m}\T{\Scene_m}{\Court_i}\mathbf{X}_k,
  \qquad
  \tilde{\mathbf{u}}_k=\Kmat\mathbf{X}^{\Cam}_k,
  \qquad
  \mathbf{u}_k=\left(\frac{\tilde u_k}{\tilde w_k},
                       \frac{\tilde v_k}{\tilde w_k}\right)
  \label{eq:label_projection}
\end{equation}
を計算する。各記録はメートル系シーン位置、カメラ深度、
画素座標と正規化座標、カメラ前方の状態、画面内の状態、
描画支持のフラグを保持する。
描画支持フラグは半径 1 画素の近傍に十分な alpha と正の描画深度があることを意味し、
点単位の深度遮蔽判定ではない。

この四つの状態は区別して保持する。

\begin{enumerate}[leftmargin=1.6em]
  \item \textbf{カメラ前方}：カメラ深度が正であること。
  \item \textbf{画面内}：投影座標が画像の内側にあること。
  \item \textbf{描画支持}：半径 1 画素の近傍に十分な alpha と正の描画深度があること。
  \item \textbf{教師有効}：上の生成物の論理積。
\end{enumerate}

この層別が重要なのは、
「見えている」と「監督してよい」が同じではないからである。
コート線の画素が遮蔽されている場合、
正解として描くべきなのは画像上で見えるマーキングだけなのか、
物理的に存在するコート構造なのかで、
教師の意味が変わる。
本研究は四つの状態を保存し、
教師有効の定義をデータ側の契約として固定する。
意味付けの選択は\cref{sec:discussion} で限界とともに述べる。

\paragraph{生成される教師}
同一の変換から次の教師が生成される。

\begin{itemize}[leftmargin=1.6em]
  \item \textbf{キーポイント}：十四チャネルのヒートマップ。
  チャネルと物理点の対応は固定で、画像座標から同一視しない。
  \item \textbf{領域分割}：背景と六種類の領域。
  コート半回転で対応する領域には同じラベルを描く。
  \item \textbf{二値の線}：コート平面上の線幅付き領域を投影したマスク。
  通常線 \(7.5\) cm、ベースライン \(15\) cm を既定とする。
  \item \textbf{線の意味分類}：どの線種かまで区別するラベル。
  \item \textbf{カメラ姿勢}：カメラ中心、カメラから正準コートへの回転、
  焦点距離。
\end{itemize}

\paragraph{姿勢の生成}
公開サンプルは
\(\T{\Scene_m}{\Cam}\)、\(\T{\Scene_m}{\Court_i}\)、\(\Kmat\) を保存し、
姿勢ベクトルを重複して持たない。
補正後にモデル側のアダプタが
\begin{equation}
  \T{\Court_i}{\Cam}=
  \bigl(\T{\Scene_m}{\Court_i}\bigr)^{-1}\T{\Scene_m}{\Cam}
  \label{eq:camera_to_court}
\end{equation}
を導出し、対象側の正準化 \cref{eq:canonical_pose} を適用する。
姿勢は三並進と三回転の六自由度を持ち、
補正後の焦点距離は剛体変換とは別に運ぶ。
補正は画像だけでなく内部パラメータと教師有効性を同時に更新し、
すべての監督量を同じ補正後フレームで扱う。

一度 \cref{eq:camera_to_court} がシーンに対して受理されれば、
追加のレンダ視点には点のクリックも線の追跡も
フレーム単位の外部校正も要らない。
人手の費用は動画の取得と、
シーン単位のアライメント検証に集約される。
この償却が、視点を大量に増やす操作を現実的にする。
